% =========================================================================
% A-4 PR-1: paper skeleton. Eight sections per drafts/paper-outline.md.
% Section bodies are paragraph-level outlines with \tbd markers — NO prose
% is written at this stage except (i) the two related-work paragraphs
% (material fixed by the outline / docs/audit/2506-21728.md) and (ii) the
% model-scoping sentences (fixed verbatim by docs/ROADMAP-B.md §0).
% Rules in force (ROADMAP-B B1, paper-outline):
%   - no theorem claim outside paper/registry.yaml (checked by
%     tools/check_paper_refs.py);
%   - "forthcoming" must not support the strength of any main theorem;
%     bounded-below / gauge normalization appears only as discussion.
% =========================================================================
\documentclass[11pt]{article}

\usepackage[margin=1.1in]{geometry}
\usepackage{amsmath,amssymb,amsthm}
\usepackage{booktabs}
\usepackage{longtable}
\usepackage{etoolbox}
\usepackage{hyperref}

% ---- Registry-generated Lean-name map + appendix table (PREAMBLE input) ----
% theorem-index.tex is AUTO-GENERATED from paper/registry.yaml by
% tools/gen_paper_index.py; it defines one \leanname@<id> macro per registry
% entry and the \TheoremIndexTable macro used in the appendix.
% AUTO-GENERATED by tools/gen_paper_index.py — do not edit by hand.
% Regenerate: python3 tools/gen_paper_index.py
% Input this file in the PREAMBLE (after \usepackage{etoolbox}). It
% defines (1) one macro \csname leanname@<id>\endcsname per registry
% entry via \csdef (consumed by the \leanref macro in main.tex),
% (2) one macro \csname leanstmt@<id>\endcsname per registry entry
% whose `latex` field is filled — entries still carrying a status
% marker are skipped, so \paperthm hard-errors on them — and
% (3) \TheoremIndexTable, which typesets the appendix table:
% paper id / Lean full name / source file.
\csdef{leanname@boundary-step-unique}{CollatzFST.boundary\_step\_unique}
\csdef{leanname@dim-l2-eq-10}{CollatzFST.Flow.finrank\_span\_dFQ\_eq\_ten}
\csdef{leanname@dim-l3-eq-31}{CollatzFST.L3.finrank\_span\_dFQ96\_eq\_31}
\csdef{leanname@flow-2k0-alternating}{CollatzFST.Flow.dF\_flow\_2K0}
\csdef{leanname@flow-kirchhoff-l2}{CollatzFST.Flow.kirchhoff\_occ2}
\csdef{leanname@fstart-l3-crossblock}{CollatzFST.L3.dF96\_fstart}
\csdef{leanname@mode-bit-sum-l3}{CollatzFST.L3.occ3\_mode\_bit\_sum}
\csdef{leanname@nogo-l2-2mode-affine}{CollatzFST.TwoMode.no\_go\_2mode\_affine\_potential}
\csdef{leanname@nogo-l2-2mode-universal}{CollatzFST.TwoMode.no\_global\_odd\_2mode\_potential}
\csdef{leanname@nogo-l2-2mode-witness}{CollatzFST.TwoMode.no\_go\_2mode\_potential}
\csdef{leanname@nogo-l2-single-universal}{CollatzFST.LP.no\_global\_odd\_ranking}
\csdef{leanname@nogo-l2-single-witness}{CollatzFST.LP.no\_nonneg\_linear\_ranking}
\csdef{leanname@nogo-l3-2mode-affine}{CollatzFST.L3.no\_go\_level3\_2mode\_affine\_potential}
\csdef{leanname@nogo-l3-2mode-universal}{CollatzFST.L3.no\_global\_odd\_level3\_2mode\_potential}
\csdef{leanname@nogo-l3-2mode-witness}{CollatzFST.L3.no\_go\_level3\_2mode\_potential}
\csdef{leanname@terminal-l2}{CollatzFST.Flow.run2\_extIn\_terminal}
\csdef{leanname@terminal-l3}{CollatzFST.L3.run3\_extIn\_terminal}
\csdef{leanname@todd-eq-dropwhile}{CollatzFST.Todd\_eq\_dropWhile}
\csdef{leanname@transducer-soundness}{CollatzFST.ofDigits\_transduce}
\csdef{leanname@transducer-split}{CollatzFST.transduce\_split}
\csdef{leanname@valuation-alt-prefix}{CollatzFST.padicValNat\_eq\_altPrefixLen}
\csdef{leanstmt@boundary-step-unique}{For every $x \in \Nat$, the number of steps $((c, P, p), b)$ of
$\microTraceTwo((1, K, 0), \extIn(x))$ with $P = K$ and
$\outBit(c, b) = 1$ equals $1$.}
\csdef{leanstmt@flow-2k0-alternating}{For every $x \in \Nat$, $\dF(x)_3 = \dF(x)_4 + \dF(x)_5$.}
\csdef{leanstmt@flow-kirchhoff-l2}{For every state $g$, every state $s \in \Seight$, and every word $w$
over $\{0,1\}$,
$\sum_{e \in \inEdges(g)} \occtwo(s, w, e) + \iv{s = g}
= \occtwo(s, w, (g, 0)) + \occtwo(s, w, (g, 1))
+ \iv{\runtwo(s, w) = g}$.}
\csdef{leanstmt@terminal-l2}{For every $x \in \Nat$,
$\runtwo((1, K, 0), \extIn(x)) = (0, S, 0)$ or
$\runtwo((1, K, 0), \extIn(x)) = (0, S, 1)$.}
\csdef{leanstmt@terminal-l3}{For every $x \in \Nat$,
$\runthree((1, K, 0, 0), \extIn(x)) = (0, S, 0, 1)$ or
$\runthree((1, K, 0, 0), \extIn(x)) = (0, S, 1, 0)$.}
\csdef{leanstmt@todd-eq-dropwhile}{For every $x \in \Nat$, $\ofDigitsTwo(\dropZeros(\transduce(x))) = \Todd(x)$.}
\csdef{leanstmt@transducer-soundness}{For every $x \in \Nat$, $\ofDigitsTwo(\transduce(x)) = 3x + 1$.}
\csdef{leanstmt@transducer-split}{For every $x \in \Nat$,
$\transduce(x) = \zeros{\vtwo(3x+1)} \listcat \dropZeros(\transduce(x))$.}
\csdef{leanstmt@valuation-alt-prefix}{For every $x \in \Nat$, $\vtwo(3x + 1) = \altPrefixLen(x)$.}
\newcommand{\TheoremIndexTable}{%
\begin{longtable}{@{}p{0.32\textwidth}p{0.63\textwidth}@{}}
\toprule
Paper id & Lean theorem / source file \\
\midrule
\endhead
\texttt{boundary-step-unique} & \texttt{CollatzFST.boundary\_step\_unique} \newline {\small \texttt{Lean4RealConstruction/Core/Collatz\_FST\_Level2.lean}} \\
\texttt{dim-l2-eq-10} & \texttt{CollatzFST.Flow.finrank\_span\_dFQ\_eq\_ten} \newline {\small \texttt{Lean4RealConstruction/ProjectA/Collatz\_FST\_DimLower.lean}} \\
\texttt{dim-l3-eq-31} & \texttt{CollatzFST.L3.finrank\_span\_dFQ96\_eq\_31} \newline {\small \texttt{Lean4RealConstruction/ProjectA/Collatz\_FST\_L3\_DimLower.lean}} \\
\texttt{flow-2k0-alternating} & \texttt{CollatzFST.Flow.dF\_flow\_2K0} \newline {\small \texttt{Lean4RealConstruction/ProjectA/Collatz\_FST\_FlowDelta.lean}} \\
\texttt{flow-kirchhoff-l2} & \texttt{CollatzFST.Flow.kirchhoff\_occ2} \newline {\small \texttt{Lean4RealConstruction/ProjectA/Collatz\_FST\_Flow.lean}} \\
\texttt{fstart-l3-crossblock} & \texttt{CollatzFST.L3.dF96\_fstart} \newline {\small \texttt{Lean4RealConstruction/ProjectA/Collatz\_FST\_L3\_Delta.lean}} \\
\texttt{mode-bit-sum-l3} & \texttt{CollatzFST.L3.occ3\_mode\_bit\_sum} \newline {\small \texttt{Lean4RealConstruction/ProjectA/Collatz\_FST\_L3\_Flow.lean}} \\
\texttt{nogo-l2-2mode-affine} & \texttt{CollatzFST.TwoMode.no\_go\_2mode\_affine\_potential} \newline {\small \texttt{Lean4RealConstruction/ProjectA/Collatz\_FST\_2Mode\_NoGo.lean}} \\
\texttt{nogo-l2-2mode-universal} & \texttt{CollatzFST.TwoMode.no\_global\_odd\_2mode\_potential} \newline {\small \texttt{Lean4RealConstruction/ProjectA/Collatz\_FST\_2Mode\_NoGo.lean}} \\
\texttt{nogo-l2-2mode-witness} & \texttt{CollatzFST.TwoMode.no\_go\_2mode\_potential} \newline {\small \texttt{Lean4RealConstruction/ProjectA/Collatz\_FST\_2Mode\_NoGo.lean}} \\
\texttt{nogo-l2-single-universal} & \texttt{CollatzFST.LP.no\_global\_odd\_ranking} \newline {\small \texttt{Lean4RealConstruction/ProjectA/Collatz\_FST\_NoLinearRanking.lean}} \\
\texttt{nogo-l2-single-witness} & \texttt{CollatzFST.LP.no\_nonneg\_linear\_ranking} \newline {\small \texttt{Lean4RealConstruction/ProjectA/Collatz\_FST\_NoLinearRanking.lean}} \\
\texttt{nogo-l3-2mode-affine} & \texttt{CollatzFST.L3.no\_go\_level3\_2mode\_affine\_potential} \newline {\small \texttt{Lean4RealConstruction/ProjectA/Collatz\_FST\_L3\_2Mode\_NoGo.lean}} \\
\texttt{nogo-l3-2mode-universal} & \texttt{CollatzFST.L3.no\_global\_odd\_level3\_2mode\_potential} \newline {\small \texttt{Lean4RealConstruction/ProjectA/Collatz\_FST\_L3\_2Mode\_NoGo.lean}} \\
\texttt{nogo-l3-2mode-witness} & \texttt{CollatzFST.L3.no\_go\_level3\_2mode\_potential} \newline {\small \texttt{Lean4RealConstruction/ProjectA/Collatz\_FST\_L3\_2Mode\_NoGo.lean}} \\
\texttt{terminal-l2} & \texttt{CollatzFST.Flow.run2\_extIn\_terminal} \newline {\small \texttt{Lean4RealConstruction/ProjectA/Collatz\_FST\_Flow.lean}} \\
\texttt{terminal-l3} & \texttt{CollatzFST.L3.run3\_extIn\_terminal} \newline {\small \texttt{Lean4RealConstruction/ProjectA/Collatz\_FST\_L3\_Flow.lean}} \\
\texttt{todd-eq-dropwhile} & \texttt{CollatzFST.Todd\_eq\_dropWhile} \newline {\small \texttt{Lean4RealConstruction/Core/Collatz\_FST\_Statements.lean}} \\
\texttt{transducer-soundness} & \texttt{CollatzFST.ofDigits\_transduce} \newline {\small \texttt{Lean4RealConstruction/Core/Collatz\_FST\_Statements.lean}} \\
\texttt{transducer-split} & \texttt{CollatzFST.transduce\_split} \newline {\small \texttt{Lean4RealConstruction/Core/Collatz\_FST\_Statements.lean}} \\
\texttt{valuation-alt-prefix} & \texttt{CollatzFST.padicValNat\_eq\_altPrefixLen} \newline {\small \texttt{Lean4RealConstruction/Core/Collatz\_FST\_Ext.lean}} \\
\bottomrule
\end{longtable}}


% ---- \leanref: cite a registry entry by its stable paper id ----
% Typesets the Lean full name and points (dagger link) to the theorem index
% in Appendix A. Unknown ids raise a hard TeX error — a second safety layer
% on top of tools/check_paper_refs.py.
\newcommand{\leanref}[1]{%
  \ifcsname leanname@#1\endcsname
    \texttt{\csname leanname@#1\endcsname}%
    \textsuperscript{\hyperref[app:theorem-index]{$\dagger$}}%
  \else
    \errmessage{leanref: unknown registry id '#1'}%
  \fi}

% ---- skeleton-stage marker for text to be written in later PRs ----
\newcommand{\tbd}[1]{\emph{[TO WRITE: #1]}}

\title{Machine-Checked Expressivity No-Go Theorems for Finite-State
  Additive Ranking Templates in Accelerated Collatz Dynamics%
  \thanks{\tbd{funding / artifact DOI footnote}}}
\author{\tbd{author list}}
\date{\tbd{date}}

\begin{document}
\maketitle

\begin{abstract}
\tbd{abstract; positioning sentence fixed by the outline: machine-checked
expressivity no-go theorems for finite-state additive ranking templates in
accelerated Collatz dynamics. Summarize the four contributions of
Section~\ref{sec:intro} and the explicit non-claims.}
\end{abstract}

% =========================================================================
\section{Introduction}\label{sec:intro}
% =========================================================================

% Outline (four contributions, in this order):
%   (1) a formalized binary carry transducer implementing the accelerated
%       Collatz map exactly (Section 2);
%   (2) a flow-geometric interpretation: occupation vectors are path flows
%       on a finite graph, and Kirchhoff conservation unifies the local
%       counting relations (Section 3);
%   (3) exact dimension theorems: the difference-feature spaces have
%       dimension exactly 10 (Level 2) and exactly 31 (Level 3 two-mode)
%       (Section 4);
%   (4) machine-checked Farkas no-go certificates for three template
%       classes, with universal quantification (x > 1) and affine
%       strengthening (Section 5).
\tbd{contribution paragraphs (1)--(4)}

% Explicit non-claims — fixed by the outline; keep both:
%   - this paper does NOT prove the Collatz conjecture, and
%   - it does NOT establish the impossibility of all finite-state methods;
%     the scope boundary is made precise in Section 6.
\tbd{non-claims paragraph}

\paragraph{Terminology.}
% Glossary fixed at skeleton stage (paper-outline.md); definitions locked,
% wording of surrounding prose free in later PRs.
\begin{itemize}
  \item \textbf{occupation vector} --- the vector of per-edge visit counts
    accumulated by one run of the transducer over an input word.
  \item \textbf{valuation-parity mode} --- the binary mode
    $m(x)$ read off the occupation feature that records the parity-relevant
    exit of the K region; the mode splits the two-register templates.
  \item \textbf{additive ranking template} --- a potential of the form
    $V(x) = \theta^{\mathsf T} F(x)$ (single mode) or
    $V(x) = \beta_{m(x)} + \theta_{m(x)}^{\mathsf T} F(x)$ (two-mode
    affine), evaluated on occupation features $F$.
  \item \textbf{accelerated Collatz map} --- $T_{\mathrm{odd}}(x) =
    (3x+1)/2^{\nu_2(3x+1)}$ on odd $x$.
  \item \textbf{Farkas certificate} --- a nonnegative rational combination
    $\lambda$ of witness rows whose aggregate annihilates all admissible
    descent directions, refuting strict descent.
  \item \textbf{mode-flow balance} --- the certificate identity
    $\sum_i \lambda_i (e_{m(y_i)} - e_{m(x_i)}) = 0$, which makes the
    per-mode intercepts $\beta_m$ free in the affine strengthening.
  \item \textbf{trimmed product} --- the useful part (reachable and
    co-reachable) of a product of a cost automaton with a language
    acceptor; used only in the discussion of Section~\ref{sec:scoping}.
\end{itemize}

% =========================================================================
\section{Transducer Semantics}\label{sec:transducer}
% =========================================================================

% Outline: the binary carry transducer realizes the accelerated map
% exactly; soundness is the registry pair
%   - \leanref of transducer-soundness : ofDigits 2 (transduce x) = 3x + 1
%   - \leanref of transducer-split     : output = v2(3x+1) zeros ++ trimmed
% plus the two semantic completions:
%   - \leanref of todd-eq-dropwhile    : trimmed value = T_odd(x)
%   - \leanref of valuation-alt-prefix : v2(3x+1) = alternating prefix len
\tbd{transducer construction paragraph}

\tbd{soundness statement paragraph, citing
\leanref{transducer-soundness} and \leanref{transducer-split}}

\tbd{accelerated-map semantics paragraph, citing
\leanref{todd-eq-dropwhile} and \leanref{valuation-alt-prefix}}

% =========================================================================
\section{Flow Conservation}\label{sec:flow}
% =========================================================================

% Outline: occupation vectors are path flows on a finite graph; Kirchhoff
% conservation unifies the local counting relations.
%   - conservation functional (feature layer): \leanref of flow-kirchhoff-l2
%   - terminal-state theorems: \leanref of terminal-l2, terminal-l3
\tbd{path-flow interpretation paragraph, citing \leanref{flow-kirchhoff-l2}}

\tbd{terminal-state paragraph, citing \leanref{terminal-l2} and
\leanref{terminal-l3}}

% MANDATORY (outline): the two structural identifications.
%   (i) the K-region alternating count e4 = e5 + e6 IS the Kirchhoff
%       conservation law at state (2,K,0) — not an independent lemma:
%       \leanref of flow-2k0-alternating;
%  (ii) boundary-step uniqueness is absorbed by the conservation row
%       space — the negative conclusion ("no separate proof needed")
%       upgraded to a positive provenance (named law, named state):
%       \leanref of boundary-step-unique.
\tbd{structural identification (i), citing \leanref{flow-2k0-alternating}}

\tbd{structural identification (ii), citing \leanref{boundary-step-unique}}

% =========================================================================
\section{Exact Geometry}\label{sec:geometry}
% =========================================================================

% Outline: exact dimensions of the difference-feature spans.
%   - Level 2: dim = 10 (in 18): \leanref of dim-l2-eq-10
%   - Level 3 two-mode: dim = 31 (in 96): \leanref of dim-l3-eq-31
\tbd{dimension theorems paragraph, citing \leanref{dim-l2-eq-10} and
\leanref{dim-l3-eq-31}}

% Cramer law (MUST be labeled a tools-layer computational result, verified
% by tools/certificates.py --cramer in CI; it is NOT a Lean theorem):
%   lambda = adjugate row of the annihilated witness submatrix / gcd;
%   certificate integers = maximal minors; in the single-mode case the
%   scale 31 IS det(A_free) — the same 31 as the lower-bound determinant.
%   Guard against the known mis-inference: sum(lambda) (1024 / 7826 / 31746)
%   is NOT any of these determinants.
\tbd{Cramer-law paragraph (tools-layer labeling as above)}

% Mode splitting breaks absorption — the three-row table (fixed by
% STATUS.md / outline):
%   single mode (48-dim): [16]+[33] holds but is absorbed by the family;
%   two-mode same-block: not even a relation (residue +-1 on mode flips);
%   two-mode cross-block: holds and is the UNIQUE gap (65th functional).
% Registry anchors: \leanref of mode-bit-sum-l3 (source identity),
% \leanref of fstart-l3-crossblock (the 65th functional as landed in Lean).
\tbd{mode-splitting three-row table and discussion, citing
\leanref{mode-bit-sum-l3} and \leanref{fstart-l3-crossblock}}

% =========================================================================
\section{No-Go Certificates}\label{sec:nogo}
% =========================================================================

% Outline: three templates, each as (finite-witness version, universal
% version over odd x > 1); affine strengthening for the two-mode templates.
%   L2 single:  \leanref of nogo-l2-single-witness / nogo-l2-single-universal
%   L2 two-mode: \leanref of nogo-l2-2mode-witness / nogo-l2-2mode-universal
%                / nogo-l2-2mode-affine
%   L3 two-mode: \leanref of nogo-l3-2mode-witness / nogo-l3-2mode-universal
%                / nogo-l3-2mode-affine
\tbd{Level 2 single-mode paragraph, citing \leanref{nogo-l2-single-witness}
and \leanref{nogo-l2-single-universal}; the universal quantifier excludes
the trivial witness $x = 1$}

\tbd{Level 2 two-mode paragraph, citing \leanref{nogo-l2-2mode-witness},
\leanref{nogo-l2-2mode-universal}, and the affine strengthening
\leanref{nogo-l2-2mode-affine}}

\tbd{Level 3 two-mode paragraph, citing \leanref{nogo-l3-2mode-witness},
\leanref{nogo-l3-2mode-universal}, and \leanref{nogo-l3-2mode-affine}}

% Mode-flow balance makes the intercepts beta_m free: the certificate
% identity sum lambda (e_m(y) - e_m(x)) = 0 (verified by
% tools/certificates.py in CI) cancels the beta contributions in the
% Farkas aggregate, so the affine class is refuted by the SAME lambda.
\tbd{mode-flow balance paragraph}

% Certificate methodology (outline): search in floating point -> exact
% rational reconstruction -> Lean decide/ring closure; anchors and CI
% enforcement (three anchor scripts + generator reproducibility).
\tbd{methodology paragraph}

% =========================================================================
\section{Model Scoping}\label{sec:scoping}
% =========================================================================

% The following classification sentences are fixed VERBATIM (English
% rendering of docs/ROADMAP-B.md §0); later PRs may add surrounding prose
% but must not weaken or restate these claims:
The single-mode no-go theorems fall in \textbf{Class A} (single-accumulator
additive rankings): $V(w) = \alpha(q_0) + \sum_i \omega(q_i, b_i) +
\beta(q_f)$.

The two-mode no-go theorems are \emph{not} Class A: in the template
$V(x) = \theta_{m(x)}^{\mathsf T} F(x)$, the K-region prefix weights depend
on the not-yet-determined mode, so the template is not any one-accumulator
edge cost. The correct classification is a \textbf{2-register copyless CRA
with final selection}: two accumulators run $\theta_0^{\mathsf T} F$ and
$\theta_1^{\mathsf T} F$ in parallel, and the mode-identifying bit selects
one of them at the final state --- the minimal edge of Class B.

Consequently, even if the universal Class-A conjecture holds, the
conclusion is the impossibility of one-accumulator additive rankings only
--- not of finite-state methods as a whole.

% Discussion-only outlook (rule: no "forthcoming" supporting main-theorem
% strength; bounded-below / gauge normalization is future work, stated as
% such). Evidence for the boundary refinement may cite the B1.5 data
% points WITH trust labels:
%   - terminal-flow imbalance of the existing certificates:
%     +-428 (Level 2, W12) and +-753 (Level 3, W20) — exact integer
%     recomputation (tools/b15_terminal_balance.py); label: [exact integer];
%   - two-balance Farkas probe (mode balance + terminal balance
%     simultaneously): feasible at both levels — floating-point LP probe
%     (HiGHS, odd x in 3..3999); label: [floating-point probe].
\tbd{discussion paragraph on gauge normalization as future work, citing the
B1.5 data points with the two trust labels above}

% =========================================================================
\section{Related Work}\label{sec:related}
% =========================================================================

% Paragraph (a) — material fixed by the outline.
Yolcu, Aaronson, and Heule (arXiv:2105.14697) establish representational
boundaries for rewriting-based matrix interpretations applied to
Collatz-like systems; the present paper establishes boundaries for
occupation-vector additive templates over a carry transducer. The two
settings differ in both certificate type and representation, but carry the
same message: for automated approaches to Collatz-style termination, the
choice of certificate class and representation imposes substantive
expressivity limits.
\tbd{one connecting sentence positioning our certificate class relative to
matrix interpretations, once Section~\ref{sec:nogo} prose is fixed}

% Paragraph (b) — VERBATIM English translation of the neutral wording in
% docs/audit/2506-21728.md («與本專案的關係» block); per the outline: no
% expansion, no point-by-point rebuttal.
That work (arXiv:2506.21728) proposes a finite-state symbolic encoding
together with a block drift; however, in its publicly available version,
the local absorption argument and the pointwise descent claim do not yet
suffice to establish the claimed global convergence; moreover, its
potential involves $\log n$ and blockwise global quantities, and thus does
not belong to the rational finite-state additive class studied in this
paper.

% =========================================================================
\section{Artifact}\label{sec:artifact}
% =========================================================================

% Outline: the artifact story — each item one short paragraph or a table:
%   - pinned toolchain (leanprover/lean4:v4.28.0-rc1) and pinned mathlib
%     revision; CI guard forbids unauthorized pin changes;
%   - zero sorry / no native_decide / no axiom beyond the kernel trio;
%     trust base of the audited theorems: propext, Classical.choice,
%     Quot.sound (Audit.lean, #print axioms);
%   - import-boundary enforcement (scripts/check_boundaries.py in CI);
%   - three anchor scripts (certificates.py, a3_functionals.py,
%     l3_recon.py) recompute every hard-coded number on every push;
%     generator reproducibility: gen_l3dim.py and gen_paper_index.py
%     outputs must match the repository byte-for-byte;
%   - the registry/PaperIndex mechanism itself: paper citations resolve
%     through paper/registry.yaml; each cited statement is restated in
%     PaperIndex.lean and kernel-checked against the original theorem, so
%     a drifting statement fails the build, and \leanref consistency is
%     enforced by tools/check_paper_refs.py. The paper thus documents its
%     own anti-drift mechanism as part of the artifact.
\tbd{artifact section body per the itemized outline above}

% =========================================================================
\appendix
\section{Theorem Index}\label{app:theorem-index}
% =========================================================================

Each \texttt{Paper.<id>} restatement is kernel-checked against the original
theorem; the table below is generated from \texttt{paper/registry.yaml}.

\TheoremIndexTable

\end{document}
